\section{Summary}
Electrical-engineering PhD researcher who bridges computational materials physics, GPU/HPC scientific software, scientific ML, and embedded engineering. Select this broad version only when a role genuinely crosses those domains.
\section{Technical Skills}
\SkillLine{Scientific compute}{Python, C++, CUDA (cuBLAS, cuTENSOR), ROCm (rocBLAS, rocFFT), MPI, OpenMP, PETSc, SLEPc, DOLFINx, JAX, PyTorch Geometric}
\SkillLine{Systems}{Frontier, Perlmutter, Docker, Kubernetes, Linux, CMake, ASP.NET, Microsoft SQL Server}
\SkillLine{Hardware}{Verilog, Lattice FPGA, PIC32, dsPIC, STM32, C8051, board debugging, OpenGL}
\section{Experience}
\Role{PhD Researcher}{Yale University}{Electrical Engineering}{2021--present}
\begin{itemize}
\item Translate many-body and first-principles physics into distributed Python/C++ workflows using MPI, OpenMP, PETSc/SLEPc, CUDA/ROCm, and national-lab GPU platforms.
\item Build DeePH-style computational graphs and scientific-ML models with equivariant GNNs and automatic differentiation for LiF and double-perovskite optical-property calculations.
\item Combine simulation research with semiconductor and photonics lab experience; publish and present self-trapped-exciton research.
\end{itemize}
\Role{Software and Embedded Systems Engineer}{Probe Technologies / Weatherford Wireline Services}{Houston, TX}{2019--2021}
\begin{itemize}
\item Shipped CUDA/OpenCL and OpenGL applications for acoustic-waveform processing and diagnostic visualization.
\item Delivered C\#, ASP.NET, and SQL Server systems for remote engineering-tool testing and sensor-calibration data.
\item Implemented firmware/FPGA functionality and supported board-level integration, debugging, and field validation.
\end{itemize}
\section{Selected Projects}
\Project{Materials ML}{Equivariant GNN and automatic-differentiation workflows for optical-property calculations.}
\Project{Integrated photonics}{Meep FDTD ring-resonator filter simulation, later fabricated and optically tested.}
\Project{Finite elements}{DOLFINx/PETSc heat, electrostatic, and flow simulations.}
\section{Education}
\Role{PhD, Electrical Engineering}{Yale University}{}{2021--present}
\Role{MPhil and MS, Electrical Engineering}{Yale University}{}{2023}
\Role{MEng and BS, Electrical and Computer Engineering}{Cornell University}{}{2013--2018}
\ifcv
\section{Research Record}
Co-author, \textit{Journal of the American Chemical Society}, 2025; APS presentations, 2024--2026.\\
\fi
